\section{Conclusion}\label{sec:conclusion}

\emph{Europe 2031} calls ASML Europe's last card. The model in this paper suggests that the card is real but that its value is set by a quantity Europe has never held in abundance: a credible alternative to Atlantic alignment. A chokepoint gives its holder leverage over the patron that depends on it only when the rival's credible offer exceeds the costs the patron can credibly impose, plus the value the holder itself places on alignment. For a Europe that depends on Washington for its defence and for frontier AI, that condition has not held since the Netherlands first let a licence for an EUV sale to China lapse without renewal in 2019. The patron's dependence on the chokepoint enters only through sanctions that would hurt the patron, the holder's agenda power, the risk of breakdown, and the holder's capacity to outlast the patron in a standoff over servicing.

The analysis connects to a broader debate about weaponized interdependence. Control of a hub is a source of power only when domestic institutions let the hub's host use it and when the host has somewhere else to go. ASML fails the first condition because its own supply chain runs through American jurisdiction, and Europe fails the second because its security does. Where those conditions change, the model predicts that the patron will try to secure control of the hub before its own threats lose their bite. In the scenario that moment comes in March 2031. In the record, the December 2024 exemption, which Washington can withdraw at will, and the 2026 fight over servicing restrictions show the United States keeping its jurisdiction in reserve and seeking to extend it to servicing while The Hague still has no exit.
